\DocumentMetadata{
  pdfversion=2.0,
  pdfstandard=ua-2,
  lang=en-US,
  tagging=on,
  tagging-setup={math/setup=mathml-SE,tabsorder=structure}
}
\documentclass[11pt,letterpaper]{article}
\usepackage[margin=0.68in,includefoot]{geometry}
\usepackage{newcomputermodern}
\usepackage{amsmath,amscd,mathtools}
\usepackage{tikz,adjustbox}
\providecommand{\square}{\mdlgwhtsquare}
\usepackage[hidelinks]{hyperref}
\hypersetup{
  pdftitle={Analysis PhD Exam, May 2018},
  pdfauthor={University of Florida Department of Mathematics},
  pdfsubject={Graduate mathematics examination},
  pdfdisplaydoctitle=true,
  bookmarksopen=true
}
\setcounter{secnumdepth}{0}
\setlength{\parindent}{0pt}
\setlength{\parskip}{0.28em}
\setlength{\emergencystretch}{3em}
\setlength{\leftmargini}{2.15em}
\setlength{\leftmarginii}{2.65em}
\setlength{\leftmarginiii}{3em}
\renewcommand{\labelenumi}{\textbf{\theenumi.}}
\pagestyle{empty}
\raggedbottom
\allowdisplaybreaks
\definecolor{ProblemConcern}{HTML}{7A1F1F}
\newenvironment{problemconcern}{%
  \par\tagstructbegin{tag=Aside}\begingroup\color{ProblemConcern}
}{%
  \par\endgroup\tagstructend\nopagebreak[4]\vspace{0.12em}
}
\newenvironment{examproblems}[1]{%
  \begin{enumerate}%
  \setcounter{enumi}{#1}%
  \setlength{\topsep}{0.35em}%
  \setlength{\partopsep}{0pt}%
  \setlength{\itemsep}{0.48em}%
  \setlength{\parsep}{0.18em}%
}{\end{enumerate}}
\newenvironment{examsubparts}{%
  \begin{enumerate}%
  \setlength{\topsep}{0.18em}%
  \setlength{\partopsep}{0pt}%
  \setlength{\itemsep}{0.16em}%
  \setlength{\parsep}{0.1em}%
}{\end{enumerate}}
\newenvironment{examinstructions}{%
  \par\begingroup\itshape
}{%
  \par\endgroup\vspace{0.18em}
}
\makeatletter
\renewcommand\section{\@startsection{section}{1}{0pt}%
  {-1.25ex plus -.25ex minus -.1ex}%
  {0.55ex plus .1ex}%
  {\normalfont\large\bfseries}}
\renewcommand{\@maketitle}{%
  \newpage\null\vskip -1em%
  {\footnotesize\bfseries
    \href{https://www.ufl.edu/}{University of Florida}, \href{https://math.ufl.edu/}{Department of Mathematics}\par}%
  \vskip -0.5em%
  \tagstructbegin{tag=Title}%
    {\LARGE\bfseries\href{https://gma.math.ufl.edu/exams/analysis-phd/}{Analysis PhD Exam}, May 2018\par}%
  \tagstructend%
  \vskip 0.25em%
  \tagmcbegin{artifact}%
    \hrule height 1pt%
  \tagmcend%
  \vskip 1em%
}
\makeatother
\title{Analysis PhD Exam, May 2018}
\author{}
\date{}
\begin{document}
\maketitle
\thispagestyle{empty}
\begin{examinstructions}
Do six of eight. Write solutions in a neat and logical fashion, giving complete reasons for all steps.
\end{examinstructions}
\begin{examproblems}{0}
\item
Suppose \((X,\mathcal M)\) and \((Y,\mathcal N)\) are measurable spaces, \(\mathcal E\subset\mathcal N\) and \(f:X\to Y\). Show, if the~\(\sigma\)-algebra generated by \(\mathcal E\) contains \(\mathcal N\) and \(f^{-1}(E)\in\mathcal M\) for all \(E\in\mathcal E\), then~\(f\) is \(\mathcal M\)-\(\mathcal N\) measurable.
\item
Do two. In each case, give a brief justification for your answer.
\begin{examsubparts}
\item[\textnormal{(a)}]
Give an example, if possible, of a measurable set \(C\subset\mathbb R\) of positive Lebesgue measure that contains no (non-trivial) interval.
\item[\textnormal{(b)}]
Give an example, if possible, of a measure space \((X,\mathcal M,\mu)\) and a sequence of measurable functions \(f_n:X\to\mathbb C\) that converge in measure, but not pointwise a.e.
\item[\textnormal{(c)}]
Give an example, if possible, of measure spaces \((X,\mathcal M,\mu)\) and \((Y,\mathcal N,\nu)\) and a function \(f:X\times Y\to[0,\infty)\) that is measurable with respect to the product measure~\(\sigma\)-algebra \(\mathcal M\otimes\mathcal N\), but for which 
\[
\int_X\int_Y f(x,y)\,d\nu\,d\mu\ne\int_Y\int_X f(x,y)\,d\mu\,d\nu.
\]
\end{examsubparts}
\item
\begin{problemconcern}
\textbf{Warning: Suspected error.} The source assigns~\(\rho\) the codomain \([0,\infty)\), although for an arbitrary \(f\in L^1(\mu)\), the values \(\int_E f\,d\mu\) need not be nonnegative. The intended repair is not uniquely determined: the codomain may be misstated, or a nonnegativity assumption on~\(f\) may be missing. The source reading has been retained.
\end{problemconcern}
Let \((X,\mathcal M,\mu)\) be a~\(\sigma\)-finite measure space (\(\mu\) a positive measure). Suppose \(\mathcal N\) is a sub-\(\sigma\)-algebra of \(\mathcal M\) and \(\nu=\mu|_{\mathcal N}\) is~\(\sigma\)-finite. Given \(f\in L^1(\mu)\), show, by considering the mapping \(\rho:\mathcal N\to[0,\infty)\) defined by \(\rho(E)=\int_E f\,d\mu\), that there exists an \(\mathcal N\)-measurable~\(g\) such that \(g\in L^1(\nu)\) and 
\[
\int_E f\,d\mu=\int_E g\,d\nu
\]
 for all \(E\in\mathcal N\). Determine~\(g\) in the case \((X,\mathcal M,\mu)\) is the interval \([0,1]\) with Lebesgue measure and \(\mathcal N=\{\varnothing,[0,\tfrac12),[\tfrac12,1],[0,1]\}\).
\item
Show \(f:[0,\infty)\to\mathbb R\) defined by \(f(x)=(1+x)^{-2}\) is integrable (with respect to Lebesgue measure on \([0,\infty)\)). Determine 
\[
\lim_{n\to\infty}\int_0^\infty \frac{\sin(x/n)}{(1+(x/n))^n}\,dx.
\]
\item
Prove, if~\(X\) is a Banach space and~\(M\) is a finite dimensional subspace of~\(X\), then there is a closed subspace~\(N\) of~\(X\) such that \(M\cap N=(0)\) and for each \(x\in X\) there exist \(m\in M\) and \(n\in N\) such that \(x=m+n\).
\item
Let \((X,\mathcal M,\mu)\) be a measure space and suppose \(1\le q<p<\infty\). Show, if \(L^p(\mu)\subset L^q(\mu)\), then there is a constant~\(\kappa\) such that if \(E\in\mathcal M\) and \(\mu(E)<\infty\), then \(\mu(E)\le\kappa\).
\item
Suppose \(E\subset\mathbb R\) is (Lebesgue) measurable with positive measure. Show 
\[
E-E:=\{x-y:x,y\in E\}
\]
 contains an open interval about~\(0\).
\item
Recall \(c_{00}\) is the vector subspace of \(\ell^\infty(\mathbb N)\) consisting of sequences \(a=(a_n)\) that are eventually zero (meaning there is an~\(N\), depending on~\(a\), such that \(a_n=0\) for \(n>N\)). Is there a norm on \(c_{00}\) that makes \(c_{00}\) a Banach space?
\end{examproblems}
\end{document}
